\begin{textsample}{POS Dim 1 – to – Score 46.00 – to/to\_ef\_32.txt}  \label{ex:f1_pos_272}
História Este texto traz alguns questionamentos acerca da \textbf{disciplina} de História e chama a atenção para o alinhamento \textbf{que} deve ser dado aos conteúdos e aos anos \textbf{que} o estudante irá percorrer ao longo da sua trajetória escolar.
Inicialmente, destacam -se alguns aspectos elencados pela professora Janice Theodoro da Silva,³ \textbf{que}, ao comentar os principais pontos \textbf{que} a BNCC estabelece para esse componente, evidencia a condição humana como \textbf{fundamento} para o currículo de História.
Para a autora, nos anos iniciais e finais do Ensino Fundamental, o professor deve possibilitar \textbf{que} o estudante desenvolva a sua concepção do “ eu ”, do “ outro ” e do “ nós ”, e, para \textbf{que} esse amadurecimento aconteça, deve -se oportunizar exercícios em sala de aula \textbf{que} estimulem o conhecimento, o respeito e a valorização de outros universos.
Ressalta também a necessidade de \textbf{que} o estudante tenha \textbf{uma} visão crítica da comunidade em \textbf{que} \textbf{vive} e trabalha e dos riscos \textbf{que} enfrenta cotidianamente.
É consenso em quase todas as partes do mundo \textbf{que} o ensino de História desenvolva nos estudantes o pensamento crítico e a autonomia intelectual.
No entanto, o objeto de estudo da \textbf{disciplina} nos debates é \textbf{que} ela se prende a fatos e personagens glorificantes.
Nessa \textbf{ótica}, como dar visibilidade aos excluídos? \textbf{Que} caminho seguir no ensino de História para preparar os cidadãos e cidadãs? \textbf{Um} dos objetivos da Base Nacional Curricular Comum ( BNCC ) é \textbf{que} o estudante se torne “ \textbf{uma} pessoa \textbf{que} identifique os desafios da vida, tenha habilidades e competências adequadas para compreender, ser capaz de analisar e encontrar soluções adequadas e éticas frente a problemas de diferentes naturezas ”.
No 1º ano do Ensino Fundamental, o universo das crianças é totalmente voltado para o lúdico, etapa em \textbf{que} elas vão tomando consciência do seu espaço na família, na comunidade ; aprendem a identificar, fazer distinção das semelhanças e diferenças entre jogos e brincadeiras \textbf{atuais} e de outras épocas, valorizando o convívio com outras crianças.
No 2º ano, as crianças aumentam seus \textbf{horizontes} e aprendem a identificar o papel de cada \textbf{um} na comunidade através dos seus registros e das experiências \textbf{vividas}.
Começam a identificar as profissões e como as mesmas \textbf{influenciam} o meio e a comunidade.
Outra habilidade \textbf{que} eles devem desenvolver refere -se às questões da sustentabilidade.
Qual deve ser o \textbf{posicionamento} de cada \textbf{um} frente ao cuidado com o meio ambiente?
Para o 3º ano, os estudantes irão mapear os grupos \textbf{que} compõem a cidade e o município.
Entender criticamente o lugar em \textbf{que} \textbf{vivem}.
Adquirir a noção do espaço público e privado e identificá -los na comunidade ; tem como foco identificar e compreender a formação e transformação na sociedade próxima à sua realidade, como exemplo podemos citar : as praças, as ruas e os monumentos.
No 4º ano, os estudantes devem ser capazes de identificar as transformações e permanências nas trajetórias ocorridas no meio social.
Estudar a circulação de pessoas, produtos e culturas.
Identificar as questões históricas relativas às migrações.
Discutir o papel da mulher e de outros grupos \textbf{historicamente} marginalizados.
No 5º ano, os estudantes têm a oportunidade de aprofundar seus conhecimentos com relação à noção de cidadania e conhecer sua \textbf{aplicabilidade} como os princípios de respeito à diversidade e aos direitos humanos.
Para tanto, o educador apresentará elementos \textbf{que} os levem a entender a cidadania como direito \textbf{que} permite ao indivíduo participar ativamente das decisões capazes de afetar a vida em sociedade, e \textbf{que} é \textbf{preciso} respeitar a diversidade e a pluralidade de pensamentos dos indivíduos.
É especificidade dos anos iniciais terem a competência e a habilidade de identificar e constituir o eu ; separar o eu do outro ; compreender a diferença ( do eu e do outro ) e aprender a conviver com ela ; perceber a distância entre \textbf{sujeito} e objeto com foco nos diferentes olhares.
Para o Ensino Fundamental, nos anos iniciais, é importante considerar a importância de se estudar o passado a partir do presente.
É nesse momento \textbf{que} o saber histórico sustentado no pensamento crítico contribui para a construção da cidadania.
Nessa perspectiva, é importante utilizar -se de diversas metodologias \textbf{que} favoreçam a apreensão do conhecimento através de distintas linguagens, tais como : leitura de obras históricas, didáticas, paradidáticas, de divulgação acadêmicas, imagens, textos, linguagem corporal, músicas, esculturas, pinturas, fotografias, mapas, tabelas, gráficos, histórias em quadrinhos ( HQ ), vídeos, infográficos, filmes, documentários, entrevistas, etc.
Ou seja, lançar mão de diferentes olhares e instrumentos, sempre na perspectiva de se formar o conhecimento histórico necessário à reflexão dos estudantes.
Embora o estudo de alguns personagens e datas seja \textbf{imprescindível} em determinados contextos, o educador não deve deixar de vinculá -los a temas significativos para a compreensão do processo histórico do estudante.
É ainda no Ensino Fundamental \textbf{que} os estudantes iniciam o estudo da história do Brasil nos seus aspectos sociais, político, econômico e cultural desde a Colônia, Império e República.
A utilização de textos de diversos \textbf{autores} da época, pesquisa iconográfica, leitura de jornais, revistas e literaturas apropriadas para a faixa etária ( obras históricas, infanto-juvenis e paradidáticas de História do Brasil ) contribuirão para a formação de leitores com \textbf{uma} visão crítica e \textbf{que} aprofundarão seus conhecimentos na segunda fase do Ensino Fundamental e no Ensino Médio.
Portanto, o professor deve estar atento à sua programação, selecionando, dos períodos históricos, aquilo \textbf{que} é mais significativo para a compreensão do processo histórico.
Havendo condições, é importante trabalhar com os estudantes na própria comunidade, realizando o estudo do lugar para \textbf{que} eles percebam os diferentes modos de \textbf{viver} no presente e em outros tempos.
Outra provocação é tentar \textbf{aproximar} algo bem distante da realidade \textbf{vivida} por eles como, por exemplo : Como são os hábitos e comportamentos das crianças \textbf{que} \textbf{vivem} e estudam na África?
E nas aldeias?
Para os estudantes do 6º ano, \textbf{um} dos elementos fundamentais e \textbf{que} percorre toda a BNCC é o conceito de comparar, entendendo ser essa a base da compreensão das diferenças.
Como por exemplo, eu comparo os meus hábitos, a minha forma de vida, com os de outras pessoas.
Cabe, ainda, aprofundar o entendimento sobre a diversidade de formas de compreensão, medição e registro do tempo, incluindo reflexões sobre sincronias e diacronias e o sentido das cronologias.
Considerando a organização dos objetos de conhecimentos elencados na BNCC, os estudantes terão oportunidades de estudar sobre as origens da espécie humana ( identificando as hipóteses científicas sobre o assunto ), seus deslocamentos e processos de sedentarização ( descrevendo as transformações da natureza pela ação humana ao longo do tempo ).
Outra habilidade \textbf{que} contribuirá para a construção de competências importantes e compreensão de nossa identidade passa pelo reconhecimento das lógicas de organização política e social, as formas de trabalho e aspectos culturais dos povos ao longo da história : Antiguidade Oriental ( mesopotâmicos, egípcios ), ( povos pré-colombianos na América ) e Antiguidade Ocidental ( gregos e romanos ) e a passagem do mundo antigo para o medieval, considerando o Mediterrâneo como espaço de interação entre Europa, África e Oriente Médio.
O estudante do 7º ano estuda o Mundo \textbf{Moderno} com algumas características mais próximas da realidade dele.
Assuntos como Renascimento e a Expansão Europeia para além-mar destacam as lógicas das sociedades daquela época e suas interações com os povos ameríndios e africanos.
No contexto da nossa história regional, abre -se espaço para o estudo das sociedades afroamericanas e indígenas, permitindo compreender suas dinâmicas sociais, suas culturas e contribuições para a formação do povo brasileiro.
No 8º ano, o estudante se debruça sobre o Mundo \textbf{Contemporâneo}, com destaque para as Revoluções Burguesas, as Rebeliões na América, as Revoluções Industriais.
É, portanto, \textbf{oportuno} \textbf{focar} os estudos na compreensão da história do mundo \textbf{contemporâneo} na Europa e na América : Iluminismo, Revoluções Inglesas, Revolução Industrial, \textbf{crise} do sistema colonial na América e os processos de independência, o Brasil monárquico, nacionalismo, liberalismo e revoluções europeias no \textbf{século} XIX, imperialismo na África e na Ásia, darwinismo social e o discurso civilizatório, e a questão indígena nas Américas.
É importante oportunizar \textbf{que} os estudantes conheçam a realidade brasileira, a fim de aprofundar análises sobre suas estruturas sociais, políticas, econômicas e culturais.
Conhecer a situação \textbf{atual} do país a partir de noticiários televisivos, pesquisa em jornais, revistas, livros didáticos e paradidáticos, filmes históricos, entrevistas com membros de determinados segmentos sociais, de acordo com a localização e possibilidades de cada escola.
No \textbf{que} se refere aos estudantes do 9º ano, por seu turno, abrem -se espaços para se trabalhar a História do Brasil, com ênfase na trajetória da República em suas diferentes fases.
Para além das esferas políticas e econômicas, temática como a Globalização transforma a sala de aula em espaço de estudos sobre problemas \textbf{contemporâneos}, especialmente a respeito dos conflitos nacionais e as pluralidades culturais, bem como as diferentes estratégias mobilizadas para resolver esses conflitos.
Na temática do “ liberalismo ”, pode -se aprofundar sobre como seu viés político e econômico se inseriu na vida brasileira, sobretudo quando o analisamos a partir dos estudos da Revolução Francesa e da Revolução Industrial.
Esses acontecimentos permitem traçar paralelos com movimentos sociais de diferentes naturezas \textbf{que} surgiram na segunda metade do \textbf{século} XVIII no Brasil, tais como a Inconfidência Mineira de 1798, com motivações essencialmente políticas, e a Revolução \textbf{Baiana}, de 1798, com sua conotação social.
Movimentos \textbf{que} mais tarde culminariam com a Independência do Brasil, em 1822.
Avançando nos estudos, o \textbf{discente} terá condições de compreender e avaliar a inserção do Brasil no sistema capitalista, como fornecedor de matérias-primas, o \textbf{que} o colocou numa relação de dependência frente aos países industrializados, a partir do \textbf{século} XIX.
Já no Tocantins, analisam -se as consequências do declínio da mineração para muitas cidades e o surgimento de \textbf{um} comércio fluvial, entre outros aspectos, \textbf{que} devem ser explorados pelos estudantes.
Outro momento \textbf{que} o estudante deve analisar e \textbf{que} faz parte da sua realidade é \textbf{conhecido} como revolução tecnológica e seus efeitos na sociedade, levando a globalização a \textbf{uma} dimensão jamais alcançada.
Tais aspectos provocam reflexões se considerarmos a globalização como \textbf{uma} moeda \textbf{que} apresenta duas \textbf{faces}, o \textbf{lado} positivo e o negativo. \textbf{Aqui}, o papel do educador é provocar as discussões sem \textbf{posicionamento} pronto ou dogmático a respeito desse processo.
Essas pesquisas e discussões podem \textbf{remeter} à trajetória do Brasil para a industrialização e sua dificuldade para deixar de ser essencialmente agrário, nos fins do \textbf{século} XIX.
Em 1º de maio de 2018, completaram -se 130 ( cento e trinta ) anos da abolição da escravidão no Brasil, porém, o Brasil foi o último país da América Latina a libertar seus escravos, e o Estado brasileiro não criou \textbf{um} mecanismo \textbf{que} garantisse a integração do negro na sociedade.
Ficando assim, essa parte da população às margens da sociedade provocando altos níveis de desigualdades \textbf{que} caracterizam a \textbf{atual} sociedade brasileira.
Trazendo para a realidade do estado, o professor tem a oportunidade de discutir, analisar e compreender com base na leitura de obras de \textbf{autores} \textbf{que} escreveram sobre a escravidão no antigo Norte goiano ( \textbf{atual} Tocantins ) a nódoa perversa do escravismo na região.
Nesse sentido, as obras da historiadora professora Juciene Ricarte Apolinário intituladas “ Vivências escravistas no Norte de Goiás no \textbf{século} XVIII”⁴ e “ A escravidão negra no Tocantins colonial”⁵ podem fornecer aos professores elementos \textbf{que} os auxiliarão na orientação do estudante quanto à análise das razões históricas \textbf{que} levam o estado do Tocantins a ranquear como o lugar com maior incidência de trabalho escravo⁶ no país.
E as mulheres? \textbf{Que} papel elas \textbf{desempenham} ao longo dos anos no Brasil desde a sua criação?
Quais as lutas \textbf{que} enfrentaram por espaço? \textbf{Que} mulheres são referências na história brasileira e \textbf{que} constam na nossa literatura?
E no Tocantins?
Esses questionamentos devem ser trabalhados com os estudantes nas diferentes séries, de maneira \textbf{que} eles possam acompanhar as lutas pelos direitos das mulheres e as conquistas obtidas.
Nesse contexto, podemos citar algumas personagens presentes na literatura e em outros espaços como : Iracema, Capitu, Marília de Dirceu, Luzia Homem, Maria Quitéria, Ana Nery, Chiquinha Gonzaga e outras personagens tocantinenses como Dona Raimunda quebradeira de coco, da região \textbf{conhecida} como Bico do Papagaio, Guilhermina Ribeiro da Silva, mais \textbf{conhecida} como dona Miúda, pioneira no artesanato de capim dourado, entre outras.
Outra \textbf{vertente} \textbf{que} traduz a realidade tocantinense dentro de \textbf{um} contexto histórico refere -se aos povos indígenas \textbf{que} compõem o estado e para maiores conhecimentos sugere -se \textbf{que} o professor conheça o livro “ Os Akroá e outros povos indígenas nas \textbf{fronteiras} do sertão”⁷.
Assim, a \textbf{disciplina} de História se caracteriza como \textbf{um} tipo de conhecimento científico construído continuamente, o \textbf{que} permite ao educador atuar como intermediador dos processos de aprendizagem, proporcionando aos educandos o contato com diferentes saberes através de variadas metodologias.
Para ensinar na perspectiva de \textbf{uma} História \textbf{que} se \textbf{atenta} para os anseios das gerações \textbf{atuais}, é necessária a construção de elos entre os acontecimentos anteriores e a realidade \textbf{contemporânea}, apresentando e debatendo os diferentes vieses \textbf{que} cada fato histórico carrega consigo.
Como mediador, o professor deve orientar os educandos a se tornarem críticos e questionarem o \textbf{que} lhes é ensinado.
A \textbf{interdisciplinaridade} no ensino de História é outro aspecto \textbf{que} deve ser considerado, \textbf{uma} vez \textbf{que} ela dialoga com as Artes, a Literatura, a Filosofia e a Sociologia, \textbf{contudo}, essa aproximação só fará sentido se o receptor for provocado a fazer interferências em sua realidade, e é o professor \textbf{que} poderá fazer essa ponte e levá -lo constantemente a ser provocado no processo de reflexão e buscas de respostas para os problemas sociais \textbf{que} afetem sua realidade.
Para o desenvolvimento de projetos interdisciplinares é necessária \textbf{uma} especial visão pedagógica e \textbf{vontade} \textbf{que} garanta \textbf{uma} postura interdisciplinar do professor ( NOGUEIRA, 2001 ).
Isso \textbf{implica} a compreensão de \textbf{que} a vida e qualquer parcela da realidade \textbf{que} se estude contêm o conhecimento como \textbf{um} todo e \textbf{que} para abordá -lo é \textbf{preciso} romper individualismos e cultivar o desejo de associar -se a outras áreas do conhecimento para tratá -lo de forma melhor e mais completa.
A história social também pode contribuir para a criticidade e a conscientização.
Todas essas provocações podem suscitar grandes mudanças no futuro desses cidadãos \textbf{hoje} em formação a partir da interferência deles na realidade em \textbf{que} estão inseridos.
Por fim, no \textbf{que} se refere à pesquisa, é importante \textbf{que} passe a fazer parte do currículo de forma regular e sistemática.
Demo ( 2000, p. 1 ) defende \textbf{que} a pesquisa deve fazer parte do perfil do profissional da Educação Básica. “ O interesse está voltado a fundamentar a importância da pesquisa para a educação, até o ponto de tornar a pesquisa à maneira escolar e acadêmica própria de \textbf{educar} ”.
Introduzir os estudantes à Iniciação Científica ( IC ) não pode ser \textbf{uma} iniciativa isolada de alguma \textbf{disciplina} ou \textbf{um} momento esporádico para alguma série desse nível de ensino, mas \textbf{uma} proposta constante do próprio projeto pedagógico da escola.
Como tal, assumida e desenvolvida pelos professores de todas as \textbf{disciplinas}, com o propósito de contribuir efetivamente com a construção das competências pertinentes ao perfil do estudante, em especial a \textbf{que} se \textbf{explicita} como a capacidade de intervir para solucionar situações complexas, valendo -se do \textbf{método} científico, bem como de propor e se envolver em alternativas criativas de melhoria do contexto individual e social, priorizando o foro da construção coletiva e da validação democrática.
Momentos privilegiados para o exercício das habilidades \textbf{que} contribuem para o desenvolvimento desta competência são as mostras científicas ou mostras de Ciências \textbf{que}, em algumas escolas, costumam ser regularmente realizadas.
Durante toda a sua trajetória pelo Ensino Fundamental, o estudante deve ser incentivado a descobrir os avanços, os limites, as ambiguidades e as incertezas \textbf{que} o \textbf{homem} carrega ao longo do tempo.
O componente curricular História deve fornecer ao jovem em formação os instrumentos para \textbf{que} lhe permita conhecer as diferentes \textbf{interfaces} \textbf{que} compõem a condição humana.
A BNCC foi pensada dentro desse grande pressuposto.
Fazer história é saber \textbf{que} a forma \textbf{que} se apreende o mundo \textbf{hoje} é bem diferente de como se apreendia na Idade Média.
Não tirar a concepção de tempo, pois história é tempo e espaço. ³ Professora titular aposentada de História da América da USP.
Disponível em https://www.youtube.com/watch? v=1XDvfxFfty4.
Acesso em 31/05/2019. ⁴ Esse texto está na coletânea organizada pelo professor Dr.
Odair Giraldin, A (trans)formação do Tocantins.
Goiânia : UFG, 2002. ⁵ Livro publicado em 2000. ⁶ Os dados foram divulgados após \textbf{um} levantamento feito pelo Ministério do Trabalho e Emprego ( MTE ), Comissão Pastoral do Trabalho ( CPT ) e o site da organização Repórter Brasil. ⁷ Livro \textbf{que} relata as políticas indígenas e indigenistas no norte da Capitania de Goiás, \textbf{atual} estado do Tocantins, \textbf{século} XVIII.
Goiânia : Kelpes, 2002.

% matched lemmas: aplicabilidade, aproximar, aqui, atentar, atual, autor, baiano, conhecido, contemporâneo, contudo, crise, desempenhar, discente, disciplina, educar, explicitar, face, focar, fronteira, fundamento, historicamente, hoje, homem, horizonte, implicar, imprescindível, influenciar, interdisciplinaridade, interface, lado, moderno, método, oportuno, posicionamento, preciso, que, remeter, sujeito, século, um, vertente, viver, vontade, ótica
\end{textsample}
